%=============================================================================
\section*{Funding}
Ministerio de Ciencia e Innovaci\'on de Espa\~na, project REBORN (PID~2021-126284OB-I00).

\section*{CRediT authorship contribution statement}
[Withheld for anonymous review]

\section*{Declaration of generative AI and AI-assisted technologies}
During the preparation of this work the authors used Claude (Anthropic) to
draft and revise sections of the manuscript and as a coding assistant. After
using this tool, the authors reviewed and edited the content as needed and take
full responsibility for the content of the published article. Language models
are also components of the system described here (statistic extraction and the
query interface); that use is described in Section~\ref{sec:system}.
\todo{Confirm this statement.}

\section*{Data availability}
CIS Study~3391: \url{https://www.cis.es}.
ISSP~2020 Environment~IV, ZA7650 v2.0.0: \url{https://doi.org/10.4232/1.14153}
(GESIS; registration required).
Computational outputs (priors, fused tables, synthetic populations, validation
results): \url{https://doi.org/10.5281/zenodo.20375257}.
Pipeline code: \url{https://anonymous.4open.science/r/virtual_surveys-5502}
[anonymised for review; public repository upon acceptance].

%=============================================================================
\bibliographystyle{apalike}
\begin{thebibliography}{99}

\bibitem[Abella et~al., 2022]{abella2025}
Abella, A., Llorach-Massana, P., Pereda-Ba\~nos, A., Marco-Almagro, L., Barreda-\'Angeles, M.,
Cl\`eries, L., 2022. Perception of recycled plastics for improved consumer acceptance through
self-reported and physiological measures. Sensors 22, 9226.
\url{https://doi.org/10.3390/s22239226}.

\bibitem[Ajzen, 1991]{ajzen1991}
Ajzen, I., 1991. The theory of planned behavior. Organ.\ Behav.\ Hum.\ Decis.\ Process.\ 50,
179--211. \url{https://doi.org/10.1016/0749-5978(91)90020-T}.

\bibitem[Barben et~al., 2008]{barben2008}
Barben, D., Fisher, E., Selin, C., Guston, D.H., 2008. Anticipatory governance of
nanotechnology: foresight, engagement, and integration, in: Hackett, E.J., Amsterdamska, O.,
Lynch, M., Wajcman, J.\ (Eds.), The Handbook of Science and Technology Studies, 3rd ed.
MIT Press, Cambridge, MA, pp.\ 979--1000.

\bibitem[Bas, 2004]{bas2004}
Bas, E., 2004. Megatendencias para el siglo XXI: un estudio Delfos. Fondo de Cultura
Econ\'omica, Mexico City.

\bibitem[Bas and Guillo, 2015]{bas2015}
Bas, E., Guillo, M., 2015. Participatory foresight for social innovation. FLUX-3D method
(Forward Looking User Experience), a tool for evaluating innovations. Technol.\ Forecast.\
Soc.\ Change 101, 275--290. \url{https://doi.org/10.1016/j.techfore.2015.06.016}.

\bibitem[Bossler et~al., 2026]{bossler2026}
Bossler, A., Fullana, A., Bas, E., 2026. When synthetic Delphi panels work: a graph-native
methodology for participatory foresight under structural complexity. Eur.\ J.\ Futures Res.\
14, 38. \url{https://doi.org/10.1186/s40309-026-00302-x}.

\bibitem[Cabrera et~al., 2025]{cabrera2025}
Cabrera, C.A.L., Ong, A.K.S., Diaz, J.F.T., Cahigas, M.M.L., Gumasing, M.J.J., 2025. Plastic to
apparel: an analysis of sustainable purchasing intention using a machine learning ensemble.
Humanit.\ Soc.\ Sci.\ Commun.\ 12, 822. \url{https://doi.org/10.1057/s41599-025-05205-z}.

\bibitem[Cao et~al., 2022]{cao2022}
Cao, Y., Lu, H., Zhu, C., 2022. Consumer preference for end-of-life scenarios and recycled
products in circular economy. Sustainability 14, 12129.
\url{https://doi.org/10.3390/su141912129}.

\bibitem[CIS, 2023]{cis2023}
Centro de Investigaciones Sociol\'ogicas, 2023. Estudio 3391: ISSP Medio Ambiente IV, Spain
[dataset]. \url{https://www.cis.es} (accessed 10 January 2024).

\bibitem[Collingridge, 1980]{collingridge1980}
Collingridge, D., 1980. The Social Control of Technology. Frances Pinter, London.

\bibitem[Di~Cicco et~al., 2024]{dicicco2020}
Di~Cicco, F., Loos, M., Vos, I., Gibson, R., Naidoo, Y., Gifford, S., Smeets, M.A.M., 2024. The
smell of sustainability: understanding consumers' acceptability of recycled plastic packaging
with an off-odor. Sustainability 16, 4896. \url{https://doi.org/10.3390/su16124896}.

\bibitem[European~Commission, 2022]{eu2022}
European Commission, 2022. Commission Regulation (EU) 2022/1616 of 15 September 2022 on
recycled plastic materials and articles intended to come into contact with foods, and
repealing Regulation (EC) No~282/2008. Off.\ J.\ Eur.\ Union L~243, 3--46.
\url{http://data.europa.eu/eli/reg/2022/1616/oj}.

\bibitem[European~Parliament and Council, 2025]{eu2025}
European Parliament and Council, 2025. Regulation (EU) 2025/40 of 19 December 2024 on
packaging and packaging waste, amending Regulation (EU) 2019/1020 and Directive (EU)
2019/904, and repealing Directive 94/62/EC. Off.\ J.\ Eur.\ Union L, 2025/40.
\url{http://data.europa.eu/eli/reg/2025/40/oj}.

\bibitem[Friedman et~al., 2008]{friedman2008}
Friedman, J., Hastie, T., Tibshirani, R., 2008. Sparse inverse covariance estimation with the
graphical lasso. Biostatistics 9, 432--441. \url{https://doi.org/10.1093/biostatistics/kxm045}.

\bibitem[Fuerth, 2009]{fuerth2009}
Fuerth, L.S., 2009. Foresight and anticipatory governance. Foresight 11~(4), 14--32.
\url{https://doi.org/10.1108/14636680910982412}.

\bibitem[Gelman et~al., 2013]{gelman2013}
Gelman, A., Carlin, J.B., Stern, H.S., Dunson, D.B., Vehtari, A., Rubin, D.B., 2013.
Bayesian Data Analysis, 3rd ed. CRC Press, Boca Raton, FL.

\bibitem[Geyer, 1991]{geyer1991}
Geyer, C.J., 1991. Markov chain Monte Carlo maximum likelihood, in: Keramidas, E.M.\ (Ed.),
Computing Science and Statistics: Proc.\ 23rd Symposium on the Interface. Interface Foundation,
Fairfax Station, pp.\ 156--163.

\bibitem[Goudie and Mukherjee, 2016]{goudie2016}
Goudie, R.J.B., Mukherjee, S., 2016. A Gibbs sampler for learning DAGs. J.\ Mach.\ Learn.\
Res.\ 17, 1--39.

\bibitem[Guston, 2014]{guston2014}
Guston, D.H., 2014. Understanding `anticipatory governance'. Soc.\ Stud.\ Sci.\ 44, 218--242.
\url{https://doi.org/10.1177/0306312713508669}.

\bibitem[Heckerman et~al., 1995]{heckerman1995}
Heckerman, D., Geiger, D., Chickering, D.M., 1995. Learning Bayesian networks: the combination
of knowledge and statistical data. Mach.\ Learn.\ 20, 197--243.
\url{https://doi.org/10.1007/BF00994016}.

\bibitem[Higgins and Thompson, 2002]{higgins2002}
Higgins, J.P.T., Thompson, S.G., 2002. Quantifying heterogeneity in a meta-analysis. Stat.\
Med.\ 21, 1539--1558. \url{https://doi.org/10.1002/sim.1186}.

\bibitem[Ibrahim and Chen, 2000]{ibrahim2000}
Ibrahim, J.G., Chen, M.-H., 2000. Power prior distributions for regression models. Stat.\ Sci.\
15, 46--60. \url{https://doi.org/10.1214/ss/1009212673}.

\bibitem[Ipsos, 2022]{ipsos2022}
Ipsos, 2022. Attitudes towards single-use plastics: a 28-country Global Advisor survey for
the Plastic Free Foundation, February 2022. Ipsos, Paris.
\url{https://www.ipsos.com/sites/default/files/ct/news/documents/2022-02/Ipsos_global_report-Attitudes_towards\%20SUP_AUSTRALIA.pdf}.

\bibitem[ISSP Research Group, 2023]{issp2023}
ISSP Research Group, 2023. International Social Survey Programme: Environment IV -- ISSP 2020.
GESIS, Cologne. ZA7650 Data file Version 2.0.0. \url{https://doi.org/10.4232/1.14153}.

\bibitem[Kjaerulff and Madsen, 2008]{kjaerulff2008}
Kjaerulff, U.B., Madsen, A.L., 2008. Bayesian Networks and Influence Diagrams: a Guide to
Construction and Analysis. Springer, New York.
\url{https://doi.org/10.1007/978-0-387-74101-7}.

\bibitem[Klinke and Renn, 2002]{klinke2002}
Klinke, A., Renn, O., 2002. A new approach to risk evaluation and management: risk-based,
precaution-based, and discourse-based strategies. Risk Anal.\ 22, 1071--1094.
\url{https://doi.org/10.1111/1539-6924.00274}.

\bibitem[Kuhlmann et~al., 2019]{kuhlmann2019}
Kuhlmann, S., Stegmaier, P., Konrad, K., 2019. The tentative governance of emerging science
and technology: a conceptual introduction. Res.\ Policy 48, 1091--1097.
\url{https://doi.org/10.1016/j.respol.2019.01.006}.

\bibitem[Leal~Filho et~al., 2021]{lealfilho2025}
Leal~Filho, W., Salvia, A.L., Bonoli, A., Saari, U.A., Voronova, V., Kl\~oga, M., Kumbhar, S.S.,
Olszewski, K., De~Quevedo, D.M., Barbir, J., 2021. An assessment of attitudes towards plastics
and bioplastics in Europe. Sci.\ Total Environ.\ 755, 142732.
\url{https://doi.org/10.1016/j.scitotenv.2020.142732}.

\bibitem[Lee et~al., 2012]{lee2012}
Lee, J., Lee, C.-Y., Lee, K.S., 2012. Forecasting demand for a newly introduced product
using reservation price data and Bayesian updating. Technol.\ Forecast.\ Soc.\ Change 79,
1280--1291. \url{https://doi.org/10.1016/j.techfore.2012.04.003}.

\bibitem[Liao et~al., 2020]{liao2020credible}
Liao, Z.A., Sharma, C., Cussens, J., van~Beek, P., 2020. Learning all credible Bayesian network
structures for model averaging. arXiv:2008.13618.

\bibitem[Liu et~al., 2009]{liu2009}
Liu, H., Lafferty, J., Wasserman, L., 2009. The nonparanormal: semiparametric estimation of
high dimensional undirected graphs. J.\ Mach.\ Learn.\ Res.\ 10, 2295--2328.

\bibitem[Lu et~al., 2012]{lu2012lsbn}
Lu, Y., Wang, M., Li, M., Zhu, Q., Yuan, B., 2012. LSBN: a large-scale Bayesian structure
learning framework for model averaging. arXiv:1210.5135.

\bibitem[Ma et~al., 2020]{ma2020}
Ma, C., Tschiatschek, S., Turner, R., Hern\'andez-Lobato, J.M., Zhang, C., 2020. VAEM: a deep
generative model for heterogeneous mixed type data. Adv.\ Neural Inf.\ Process.\ Syst.\ 33,
11925--11936.

\bibitem[Magnier et~al., 2019]{magnier2019}
Magnier, L., Mugge, R., Schoormans, J., 2019. Turning ocean garbage into products:
consumers' evaluations of products made of recycled ocean plastic. J.\ Clean.\ Prod.\ 215,
84--98. \url{https://doi.org/10.1016/j.jclepro.2018.12.246}.

\bibitem[Mao et~al., 2020]{mao2020}
Mao, C., Koide, R., Brem, A., Akenji, L., 2020. Technology foresight for social good:
social implications of technological innovation by 2050 from a global expert survey.
Technol.\ Forecast.\ Soc.\ Change 153, 119914.
\url{https://doi.org/10.1016/j.techfore.2020.119914}.

\bibitem[Martins et~al., 2024]{martins2024}
Martins, L.N.A., Gon\c{c}alves, F.B., Galletti, T.P., 2024. Generation and analysis of
synthetic data via Bayesian networks: a robust approach for uncertainty quantification via
Bayesian paradigm. Preprint arXiv:2402.17915.
\url{https://doi.org/10.48550/arXiv.2402.17915}.

\bibitem[Meng and Leary, 2021]{meng2021}
Meng, M.D., Leary, R.B., 2021. It might be ethical, but I won't buy it: perceived
contamination of, and disgust towards, clothing made from recycled plastic bottles.
Psychol.\ Mark.\ 38, 298--312. \url{https://doi.org/10.1002/mar.21323}.

\bibitem[Molenbuur et~al., 2025]{molenbuur2025}
Molenbuur, M., van~Eijk, M.C.P., Urbanus, J.H., Diepenmaat, H., Ragaert, K., 2025.
Perceptions and alignment on quality along the circular plastics packaging material chain.
Waste Manag.\ 200, 114758. \url{https://doi.org/10.1016/j.wasman.2025.114758}.

\bibitem[Neumann et~al., 2024]{neumann2024}
Neumann, N., Mangold, A., Schneider, J., W\"olfel, C., 2024. Driving a circular economy: nudging
and innovative strategies for promoting post-consumer recycled plastic home-appliances.
Circ.\ Econ.\ Sustain.\ 5, 1607--1629. \url{https://doi.org/10.1007/s43615-024-00487-y}.

\bibitem[Pearl, 1988]{pearl1988}
Pearl, J., 1988. Probabilistic Reasoning in Intelligent Systems: Networks of Plausible
Inference. Morgan Kaufmann, San Francisco, CA.

\bibitem[Polyportis et~al., 2022]{polyportis2022}
Polyportis, A., Mugge, R., Magnier, L., 2022. Consumer acceptance of products made from
recycled materials: a scoping review. Resour.\ Conserv.\ Recycl.\ 186, 106533.
\url{https://doi.org/10.1016/j.resconrec.2022.106533}.

\bibitem[Quay, 2010]{quay2010}
Quay, R., 2010. Anticipatory governance: a tool for climate change adaptation.
J.\ Am.\ Plan.\ Assoc.\ 76, 496--511. \url{https://doi.org/10.1080/01944363.2010.508428}.

\bibitem[Reijonen et~al., 2021]{reijonen2020}
Reijonen, H., Bellman, S., Murphy, J., Kokkonen, H., 2021. Factors related to recycling plastic
packaging in Finland's new waste management scheme. Waste Manag.\ 131, 88--97.
\url{https://doi.org/10.1016/j.wasman.2021.05.034}.

\bibitem[Roche~Cerasi et~al., 2021]{plasticircle2023}
Roche~Cerasi, I., S\'anchez, F.V., Gallardo, I., G\'orriz, M.\'A., Torrijos, P., Aliaga, C.,
Franco, J., 2021. Household plastic waste habits and attitudes: a pilot study in the city of
Valencia. Waste Manag.\ Res.\ 39, 679--689. \url{https://doi.org/10.1177/0734242X21996415}.

\bibitem[Rogge and Reichardt, 2016]{rogge2016}
Rogge, K.S., Reichardt, K., 2016. Policy mixes for sustainability transitions: an extended
concept and framework for analysis. Res.\ Policy 45, 1620--1635.
\url{https://doi.org/10.1016/j.respol.2016.04.004}.

\bibitem[Ruokamo et~al., 2022]{ruokamo2022}
Ruokamo, E., R\"ais\"anen, M., Kauppi, S., 2022. Consumer preferences for recycled plastics:
observations from a citizen survey. J.\ Clean.\ Prod.\ 379, 134720.
\url{https://doi.org/10.1016/j.jclepro.2022.134720}.

\bibitem[Stern, 2000]{stern2000}
Stern, P.C., 2000. Toward a coherent theory of environmentally significant behavior. J.\ Soc.\
Issues 56, 407--424. \url{https://doi.org/10.1111/0022-4537.00175}.

\bibitem[Stilgoe et~al., 2013]{stilgoe2013}
Stilgoe, J., Owen, R., Macnaghten, P., 2013. Developing a framework for responsible
innovation. Res.\ Policy 42, 1568--1580. \url{https://doi.org/10.1016/j.respol.2013.05.008}.

\bibitem[Walker et~al., 2013]{walker2013}
Walker, W.E., Haasnoot, M., Kwakkel, J.H., 2013. Adapt or perish: a review of planning
approaches for adaptation under deep uncertainty. Sustainability 5, 955--979.
\url{https://doi.org/10.3390/su5030955}.

\bibitem[Widiastuti et~al., 2025]{widiastuti2025}
Widiastuti, I., Azzam, Y.F., Aldiyanto, D.N., Saputra, T.W., Estriyanto, Y., 2025. Academic
community acceptance of recycled plastic products: an analysis of product evaluation and
willingness to consume. Clean.\ Waste Syst.\ 11, 100280.
\url{https://doi.org/10.1016/j.clwas.2025.100280}.

\bibitem[Xu et~al., 2019]{xu2019}
Xu, L., Skoularidou, M., Cuesta-Infante, A., Veeramachaneni, K., 2019. Modeling tabular data
using conditional GAN, in: Wallach, H., et~al.\ (Eds.), Adv.\ Neural Inf.\ Process.\ Syst.\
32. Curran Associates, Red Hook, NY, pp.\ 7335--7345.

\bibitem[Zhang et~al., 2017]{zhang2017}
Zhang, J., Cormode, G., Procopiuc, C.M., Srivastava, D., Xiao, X., 2017. PrivBayes: private
data release via Bayesian networks. ACM Trans.\ Database Syst.\ \chk{42(4), 25}.
\url{https://doi.org/10.1145/3134428}.

\bibitem[Zwicker et~al., 2020]{zwicker2020}
Zwicker, M.V., Nohlen, H.U., Dalege, J., Gruter, G.-J.M., van~Harreveld, F., 2020. Applying
an attitude network approach to consumer behaviour towards plastic. J.\ Environ.\ Psychol.\
69, 101433. \url{https://doi.org/10.1016/j.jenvp.2020.101433}.

\end{thebibliography}

%=============================================================================
\appendix
\section{External studies}
\label{app:studies}

\begin{table}[htbp]
\centering
\caption{Extraction corpus. $n$: largest survey sample used.
\emph{Entries}: extracted item-level distributions. \emph{Retained evidence}:
target outcomes for which the study supplies entries after the review rules.
Ipsos (2022) is used only as an external benchmark
and \citet{polyportis2022} only for the construct taxonomy; neither supplies
priors.}
\label{tab:studies}
\small
\begin{threeparttable}
\begin{tabular}{p{3.9cm}p{2.1cm}rrp{5.6cm}}
\toprule
\textbf{Study} & \textbf{Country}\tnote{a} & $n$ & \textbf{Entries} &
  \textbf{Retained evidence}\tnote{b} \\
\midrule
%%STUDYROWS%%
\bottomrule
\end{tabular}
\begin{tablenotes}\small
\item[a] As recorded in the extraction metadata of each source.
\item[b] Entries for CIS-observed variables were extracted but are not used
  (Section~\ref{sec:step3}). ``None'': no entry retained. The base rate and
  coefficients of Eq.~\eqref{eq:transfer} are all from \citet{ruokamo2022}.
\end{tablenotes}
\end{threeparttable}
\end{table}

\section{Sensitivity analyses}
\label{app:sensitivity}

\begin{table}[htbp]
\centering
\caption{Sensitivity of projected food-packaging acceptance to the transfer
factor $a_0$ (with $p_0 = 0.38$) and to the base rate $p_0$ (with
$a_0 = 0.80$). The population mean equals $p_0$ by calibration. The ordering of
the two profiles is preserved throughout. Source:
\texttt{transfer\_output/sensitivity\_*.csv}.}
\label{tab:a0_sensitivity}
\small
\begin{tabular}{lcccc}
\toprule
 & \textbf{Value} & P(most receptive) & P(least receptive) & \textbf{Spread} \\
\midrule
%%A0ROWS%%
\midrule
%%P0ROWS%%
\bottomrule
\end{tabular}
\end{table}


\begin{table}[htbp]
\centering
\caption{Projected food-packaging acceptance by segment under each transfer
factor $a_0$ ($p_0 = 0.38$; the population mean is fixed at $p_0$). Segment
contrasts widen with $a_0$; their ordering is unchanged. Source:
\texttt{transfer\_output/segments\_a0.csv}.}
\label{tab:seg_a0}
\small
\begin{tabular}{lrcccc}
\toprule
\textbf{Segment} & $n$ & $a_0 = 0.50$ & $0.65$ & $0.80$ & $1.00$ \\
\midrule
%%SEGA0ROWS%%
\bottomrule
\end{tabular}
\end{table}

\section{Missing values}
\label{app:missing}

\begin{table}[htbp]
\centering
\caption{Network variables with more than 10\% missing values in CIS~3391,
imputed with the modal category. %%MISSSUMMARY%%}
\label{tab:missing}
\small
\begin{tabular}{lr}
\toprule
\textbf{Variable} & \textbf{Missing (\%)} \\
\midrule
%%MISSROWS%%
\bottomrule
\end{tabular}
\end{table}

\section{Pipeline scripts}
\label{app:code}

\begin{table}[H]
\centering
\caption{Pipeline scripts and primary outputs (repository \texttt{virtual\_surveys}).\label{tab:scripts}}
\small
\begin{tabular}{lll}
\toprule
\textbf{Script} & \textbf{Section} & \textbf{Key output} \\
\midrule
\texttt{ingest\_pdf.py}           & \ref{sec:system}       & Text extraction from PDF sources \\
\texttt{harmonize\_categories.py} & \ref{sec:step2}        & \texttt{category\_mappings.json} \\
\texttt{step1.py}                 & \ref{sec:step1}        & \texttt{dag\_edges\_3391.csv} \\
\texttt{step2.py}                 & \ref{sec:step2}        & \texttt{step2\_dirichlet\_priors.json} \\
\texttt{evidence\_table.py}       & \ref{sec:step2}        & \texttt{target\_evidence.csv} \\
\texttt{step4b\_targets.py}       & \ref{sec:step4}        & \texttt{synthetic\_surveys\_final.csv} \\
\texttt{step4.py}                 & \ref{sec:step4}        & \texttt{synthetic\_surveys\_S1.csv} \\
\texttt{step6.py}                 & \ref{sec:res_fidelity} & \texttt{tvd\_per\_variable\_fixed.csv} \\
\texttt{step8.py}                 & \ref{sec:res_fidelity} & \texttt{comparison\_table.csv} \\
\texttt{step9.py}                 & \ref{sec:res_ipsos}    & \texttt{validation\_final\_S1\_vs\_ipsos.json} \\
\texttt{transfer\_model.py}       & \ref{sec:step2}, \ref{sec:spread} & \texttt{transfer\_output/*.csv} \\
\texttt{transport\_issp.py}       & \ref{sec:transport}    & \texttt{issp\_transport.json} \\
\texttt{transfer\_transport.py}   & \ref{sec:sensitivity}  & \texttt{transport\_segments.csv} \\
\texttt{measurement\_plan.py}     & \ref{sec:implications} & \texttt{measurement\_plan.csv} \\
\texttt{app\_gradio.py}           & \ref{sec:system}       & Query interface (ASAP~v6) \\
\bottomrule
\end{tabular}
\end{table}

\section{Technical details of the inference engine}
\label{app:technical}

\subsection{Structure learning: parallel-tempered MC$^3$}

We approximate the posterior over DAGs $p(G \mid X)$ using a parallel-tempered MC$^3$
variant \citep{geyer1991} inspired by the Gibbs sampler of \citet{goudie2016}, employing
multiple chains at temperatures $T \in \{1.0, 1.5, 3.0, 8.0, 25.0, 100.0\}$ with periodic
swap proposals to escape local modes. The cold chain ($T=1$) targets the actual posterior;
hot chains ($T>1$) explore a flattened version of the posterior landscape. A
sparsity-inducing prior penalises dense graphs:
\begin{equation}
p(G) \propto \prod_j \exp\!\bigl(-\gamma\,|\mathrm{pa}(X_j)|\bigr),
\end{equation}
with $\gamma$ calibrated via cross-validation. The in-degree bound is set
data-adaptively as $d^* = \lfloor \log(n / K\tau) / \log K \rfloor$, which
gives $d^* = 3$ within blocks for the CIS data; in the ensemble of graphs used
for generation the maximum in-degree is 2.

\subsection{Thematic block decomposition}

The DAG space for 90 variables is too large for a monolithic sampler: already at $d=10$
there are over $4 \times 10^{18}$ possible DAGs \citep{kjaerulff2008}. Variables are
partitioned into 15 thematic blocks of at most 12 variables each, following the structure of
the ISSP questionnaire: two environmental concern blocks, willingness to pay and acceptance,
behaviour, ecological footprint, perceived environmental danger, science and technology
attitudes, institutional trust, political issues, political values, political ideology, life
satisfaction, employment characteristics, family origin, and demographics. Demographic root
variables (SEX, BIRTH, NAT\_INC, ESTUDIOS) are co-located within the thematic blocks most
relevant to each, so that the sampler can learn dependencies between demographics and
attitudes within a block.

Cross-block dependencies are identified through a four-layer procedure. First, all variables
are mapped to a latent Gaussian space via rank-based normal scores (the nonparanormal
transformation; \citealt{liu2009}), and the graphical lasso \citep{friedman2008} estimates
the sparse precision matrix. Second, pairs with high mutual information not captured by the
lasso are added as candidates. Third, each candidate edge is validated through
bootstrap-based BDeu score comparison, retaining edges that improve local scores in at least
80\% of resamples. Fourth, a greedy top-3-per-node constraint bounds the number of
cross-block parents. No edges are imposed. The target outcomes are not nodes of the
network; they are attached through Eq.~\eqref{eq:transfer}.

\subsection{Posterior predictive generation}

Synthetic individuals are generated by ancestral sampling. For each of the $M=50$ graphs
drawn from the posterior,
\begin{equation}
X_j \sim \mathrm{Cat}\bigl(\mathbb{E}[\boldsymbol{\theta}_{j,c}]\bigr),
\end{equation}
using the posterior mean of each conditional table for computational tractability, following
the efficient variant of \citet{martins2024}. Intervals for the target outcome are computed
by Monte Carlo as described in Section~\ref{sec:step5}.

\subsection{Total variation distance}

For two categorical distributions $p$ and $q$,
\begin{equation}
\mathrm{TVD}(p,q) = \tfrac{1}{2}\sum_k |p_k - q_k|.
\end{equation}

\end{document}
